\documentclass[11pt, a4paper]{article}
\input{preamble.tex}
\switchlang{en}
\begin{document}

\title{1989 -- Numonics Manager Mouse Cordless, 2nd generation}
\date{}
\maketitle
\selectlanguage{english}

The second-generation Numonics Manager Mouse is a redesigned version of the original Manager Mouse released in 1985 by the Torrington Company, a major international bearing manufacturer headquartered in Connecticut \cite{torringtonAdd}. Since mouse production was far away from Torrington's core business, they eagerly licensed  Manager Mouse to other companies, such as Maynard Electronics \cite{admaynard} and Numonics Corporation \cite{adnumonics}. Numonics specialized in precision measuring instruments and also gained recognition as a manufacturer of CAD computer peripherals, such as digitizers and plotters \cite{numonics}. Initially, it included a Torrington-designed mouse in the range of devices under its brand, and then in 1989 it introduced an updated model of the mouse under the same name, the Numonics Manager Mouse. This updated version is not available under other manufacturers' brands, being an exclusive device, and is intended for CAD and graphics work \cite{buxtonNumonicsCardboard}, which is logical, given the company's focus. The FCC ID code present on the body also belongs to Numonics Corporation, but Torrington's involvement as the actual manufacturer of the device is still quite possible.

\begin{figure}[h]
   \centering
    \includegraphics[scale=0.9]{1989_numonics_manager_mouse_cordless_2gen/pic_30.jpg}
    \caption{Manager Mouse}
    \label{fig:NumonicsManagerMouse2Pic}
\end{figure}

Like the first generation Manager Mouse, this mouse also appeared on the market in two versions: a wired version (Manager Mouse 3S) and a wireless infrared version (Manager Mouse Cordless) \cite{italianReview}. The model was available for IBM compatible computers and for the Apple Macintosh (the latter only in a wired version) \cite{buxtonNumonicsCardboard}. Following the first Manager Mouse model, the mouse retained its unusual wheel-based way of motion detection.

The year when the mouse had appeared can be figured out using the FCC ID code, which, according to the Federal Communications Commission database, shows that the mouse was released in 1989, as well as by the postmark on the scanned version \cite{buxtonNumonicsCardboard} of the packaging.

\section*{Exterior}

The mouse body is made of white plastic. Like the first generation mouse, the case of the Manager Mouse 2nd generation has a rectangle in its base, but that's where the similarities end. The upper part of the body has a streamlined ``aerodynamic'' shape; its longitudinal cross-section is an asymmetrical bell-shaped profile that rises near the user's wrist. Three large buttons, the same color as the body, are fully integrated into its shape and have no protrusions (the middle button has a round bulge for easier tactile identification). There are no other elements on the upper side of the body, except for the name ``Manager Mouse'' painted near the left button. The front (farthest from the user) part of the body contains an infrared emitter: therefore, it is covered by a convex ``bumper'' made of glossy black plastic (the wired version of the mouse has an entirely white body). The bumper is illuminated by LEDs during information exchange, acquiring a translucent bluish tint \cite{italianReview}. In the center of the bumper (where the wired version has exit of the mouse cable) there is a hole through which the charging cable plug is inserted.

\begin{figure}[h]
    \centering
    \includegraphics{1989_numonics_manager_mouse_cordless_2gen/top_15.jpg}
    \includegraphics{1989_numonics_manager_mouse_cordless_2gen/bottom_15.jpg}
    \caption{Manager Mouse, top and bottom views}
    \label{fig:NumonicsManagerMouse2TopAndBottom}
\end{figure}

The bottom of the mouse is very similar to the Torrington Manager Mouse (fig. \ref{fig:NumonicsManagerMouse2TopAndBottom}): it features a sticker with informational data, round low-friction pads that make the mouse glide easier, as well as two small truncated cone-shaped wheels designed to register movement, instead of the traditional ball of other mechanical mice. This design comes back to Douglas Engelbart's first mice: the wheels are positioned orthogonally to each other, so that one always rotates while the other always slides during the longitudinal and lateral movements. However, compared to Engelbart's mouse, the Manager Mouse's wheels are significantly smaller, and their axes are angled approximately 15 degrees to the horizontal plane \cite{italianReview}, which should improve the registration of diagonal movements.

\section*{Size and ergonomics}

\begin{figure}[h]
    \centering
    \includegraphics[scale=0.56]{1989_numonics_manager_mouse_cordless_2gen/size_15.jpg}
    \caption{Manager Mouse on a graduated pad with a grid step of 1~cm}
    \label{fig:NumonicsManagerMouse2Size}
\end{figure}

The case is of a size typical for the second half of the 1980s (fig. \ref{fig:NumonicsManagerMouse2Size}). However, thanks to significant changes in geometry compared to the first generation of the Manager Mouse, the shape of this mouse can be considered one of the most ergonomic among the mice of the late 1980s. This applies to the wider buttons integrated into the case and the convex palm rest. Furthermore, the mouse body is symmetrical, equally comfortable for both left- and right-handed users.

The review \cite{italianReview} notes that the features of the motion detection mechanism ensure very precise positioning of the pointer, but the design of the wheels still has some drawbacks: they do not respond perfectly to diagonal movements (thus, the angled position of the axes, intended to solve this problem, does not cope with the task perfectly).

\begin{figure}[h]
    \centering
    \includegraphics[scale=0.56]{1989_numonics_manager_mouse_cordless_2gen/hand_15.jpg}
    \caption{Manager Mouse with a human hand model}
    \label{fig:NumonicsManagerMouse2Hand}
\end{figure}

\section*{Connectivity}

Both Torringon/Numonics in 1986 and even Numonics in 1989 advertised the wireless version of the Manager Mouse as the ``world's first wireless mouse'' (in reality, the first mouse with such an interface was an infrared mouse developed by Logitech in 1984 for Metaphor Computer Systems workstations).  However, like the wired Manager Mouse, the wireless version's infrared receiver connects to the computer via an RS-232 port (both Microsoft and Mouse Systems protocols are supported). The receiver cable has a 9-pin serial connector, which also contains a socket for connecting a 9-volt power supply. An adapter to a 25-pin RS-232 connector and a charging cable are also included. To charge the mouse, this cable is connects the receiver and the jack on the front of the mouse. It is emphasized in the review \cite{italianReview} that charging does not require a serial port connection, and that a full charge takes at least 15 hours, after which the device can operate for approximately 8 hours. It automatically enters standby mode after one minute of inactivity, after which it must be activated by pressing the left button.

\section*{Internal design}

The internal structure of the mouse is shown in fig. \ref{fig:NumonicsManagerMouse2Inside}. As you can see, it features an optomechanical design (unlike the mechanical encoders of the first-generation Manager Mouse), but the encoder layout remains the same: the optical interrupter disks are mounted directly to the wheel axes. The mouse's physical resolution has been doubled compared to the first-generation model, reaching 200 DPI \cite{resolution}.

\begin{figure}[h]
    \centering
    \includegraphics{1989_numonics_manager_mouse_cordless_2gen/inside_30.jpg}
    \caption{Manager Mouse disassembled}
    \label{fig:NumonicsManagerMouse2Inside}
\end{figure}

The photo also shows a pack of four rechargeable Ni-Cd batteries housed in a removable black plastic container, which significantly increases the mouse's weight. An additional printed circuit board, located on top of the mouse's main circuit board at the front of the body, is used to charge them and also provides data exchange. It contains a charging connector and two infrared emitting diodes positioned at a 60-degree angle \cite{italianReview}.

\begin{thebibliography}{9}
\bibitem {torringtonAdd} The smallest link in your desktop publishing system is the strongest (advertising) // PC Magazine, V. 6, No. 17, October  13, 1987. - p. 305.\url{https://archive.org/details/PC-Mag-1987-10-13/page/304/mode/2up}

\bibitem {admaynard} This Little Fella Means Business [adv.] // PC Magazine, Vol. 5, No. 1, January 14, 1986. P. 37. \url{https://archive.org/details/PC-Mag-1986-01-14/page/n71/mode/2up}

\bibitem {adnumonics} Don't let your mouse get hung-up on itself [adv.] // PC Magazine, Vol. 7, No. 6, March 29, 1988. P. 282 // \url{https://archive.org/details/PC-Mag-1988-03-29/page/n277/mode/2up}

\bibitem {numonics} About Numonics Corporation -- numonics.com \url{https://web.archive.org/web/19981203122830/http://www.numonics.com/}

\bibitem {buxtonNumonicsCardboard} NumonicsCordlessMouse053 -- billbuxton.com \url{https://www.billbuxton.com/inputTimelineAssets/NumonicsCordlessMouse053.pdf}

\bibitem{italianReview} MC Microcomputer, No. 87, July-August 1989. P. 68. \url{https://archive.org/details/MC_microcomputer-087/page/n67/mode/2up}

\bibitem{resolution} Tout sur les souris // Compatibles PC Magazine, No. 29, Decembre 1989. P. 19 \url{https://archive.org/details/compatibles-pc-magazine-29-decembre-1989/page/n19/mode/2up}
\end{thebibliography}
\end{document}
